% interface=en output=pdftex

\environment mstyle 

\starttext 

\TitlePage{507}{0}{\TeX WORK\\explained}{}{1}

\subject {Introduction}

The \TEXWORK\ environment is primary meant for use within 
\PRAGMA, but can be of use to others too. When we started 
using Personal computers back in the 80's, there were not so 
many text editors around. Therefore we wrote our own editing 
environment, which could handle very large documents, even 
on small machines. Also, we favored a systematic key binding
over cryptic control and alt key based bindings, especially
for frequent actions. 

When we started using \TEX, we extended this editor with 
some basic features, like syntax highlighting, use of 
symmetry and fast directory browsing. As \CONTEXT\ grow, we 
also added template support, real time spell checking, 
\CONTEXT\ project support, and help information. Apart from 
some additional features, \TEXWORK, is based on this editor, 
called \TEXEDIT. 

Where \TEXEDIT\ was written in Modula~2, and all screen 
handling was taken care of by the program itself, \TEXWORK\ 
uses \PERLTK. This enabled us to migrate from the \MSDOS\ 
operating system to \LINUX\ while maintaining our methods of 
working. 

\subject {General Structure}

Because \TEXWORK\ is modelled after \TEXEDIT, is has the
same largely symmetrical keybindings. The most common
editing actions are one||key ones. When in character mode,
the \type {delete} key deletes the next character, and
\type{insert} inserts the last deleted one. The \type
{delete}, \type {insert}, \type{begin} and \type {end}. 

\subject {Key bindings}

The function keys are used for setting the current mode. The 
first three modes speak for themselves. The page mode is 
limited to the formfeed character, represented by \type 
{\f} in many programming languages. We sometimes use this
character as a visual page separator, but for \CONTEXT\ this
character is just equal to \type{\par} (an empty line). 

\starttabulate
\NC F1  \NC enable character mode   \NC \NR             
\NC F2  \NC enable word mode        \NC \NR
\NC F3  \NC enable line mode        \NC \NR       
\NC F4  \NC enable page mode        \NC \NR
\NC F5  \NC enable paragraph mode   \NC \NR
\NC F6  \NC enable TeX/MP mode      \NC \NR
\stoptabulate      

The fifth mode concerns paragraphs, and here \type {delete} 
and \type {insert} are disabled. It is not really funny to 
see the paragraphs disappear rappidly when we keep the \type 
{delete} key pressed.

Mode~6, duplicated in~12 because this key is conveniently 
positioned at the edge of the row, concerns \TEX. Apart from
the menu actions, the keys are bound like in word mode. 

\starttabulate
\NC F7  \NC enable buffer mode      \NC \NR
\NC F8  \NC enable text mode        \NC \NR
\NC F9  \NC toggle select mode      \NC \NR
\stoptabulate      

Buffer mode is there when we want to manipulate a selected 
area. By clicking the bottom left area of the screen, right 
below the line numbers, you can change the line based buffer
(and selection) mode into a blockwise alternative. Text mode
concerns the whole text. 

You can select text with either the mouse, or with the 
cursor keys and the \type{begin} and \type{end} keys. A 
selection can be deleted with \type {delete}, in which case 
the content is moved into the selection buffer. When \type 
{insert} is pressed, the content is only copied. When 
nothing is selected, and buffermode is on, \type {insert} 
copies the buffer into the text. When the \type {control} 
key is pressed too, the clipboard is used instead of the 
buffer. 

Each mode has its own actions, shown at the bottom of the 
screen. You can walk though the actions with the tenth 
function key. An action is applied or repeated with \type 
{F11}. 

\starttabulate
\NC F10 \NC set (next) menu action  \NC \NR  
\NC F11 \NC repeat last menu action \NC \NR  
\stoptabulate      

As said, the twelfths key act just like number six. 

\starttabulate
\NC F12 \NC TeX/MP mode             \NC \NR
\stoptabulate      

The meaning of the keypad keys will be no surprise: 
             
\starttabulate
\NC Delete \NC delete current object    \NC \NR 
\NC Insert \NC re-insert current object \NC \NR 
\NC Home   \NC goto begin of object     \NC \NR 
\NC End    \NC goto end of object       \NC \NR 
\NC Prior  \NC half a screen backwards  \NC \NR 
\NC Next   \NC half a screen forward    \NC \NR 
\stoptabulate

Normally all actions work top||down, but you can revert
that direction with:

\starttabulate
\NC Control-Prior \NC work backwards \NC \NR 
\NC Control-Next  \NC work forward   \NC \NR 
\stoptabulate

The tab keys are not inserting tabs, but skip to the next
column, as derived from the previous of next non empty line. 

\starttabulate
\NC Tab       \NC align to next word in previous line \NC \NR 
\NC Shift-Tab \NC align to next word in next line     \NC \NR 
\stoptabulate

The next few bindings influence the size of the working area
and the size of the font. 

\starttabulate
\NC Alt-Left  \NC use a smaller text font \NC \NR 
\NC Alt-Right \NC use a larger text font  \NC \NR 
\NC Alt-Up    \NC minimize screen         \NC \NR 
\NC Alt-Down  \NC maximize screen         \NC \NR 
\stoptabulate

You can change the focus between the log, search and text 
windows by means of:

\starttabulate
\NC Shift-Left  \NC cycle screens backward \NC \NR 
\NC Shift-Right \NC cycle screens forward  \NC \NR 
\stoptabulate

For power users who do a lot of editing, the next keys can
be of help. If you forget those, it does not harm. 

\starttabulate
\NC Control-Up    \NC previous file in menu  \NC \NR 
\NC Control-Down  \NC next file in menu      \NC \NR 
\NC Control-Left  \NC cycle file menus left  \NC \NR 
\NC Control-Right \NC cycle file menus right \NC \NR 
\stoptabulate

When error checking is done (in \TEX/\METAPOST\ mode) or when all 
files are searched, you can follow the results with: 

\starttabulate
\NC Control-Home \NC goto end of next/prev error/found   \NC \NR 
\NC Control-End  \NC goto begin of next/prev error/found \NC \NR 
\stoptabulate

Swhitching to block selection mode can be accomplished by
clicking on the bottom left rectangle, but \type {Control-b}
will do too. Another control key concerns dislexis mode (one
of the oldest \TEXEDIT\ functions).  

\starttabulate
\NC Control-b \NC toggle block select mode (big white/green button) \NC \NR 
\NC Control-w \NC swap characters                                   \NC \NR 
\stoptabulate

We already mentioned that selected content can be moved to 
the buffer or clipboard. The clipboard is a system wide 
buffer, and can be used to move data between applications. 

\starttabulate
\NC Control-Insert \NC clipboard to text/selection to clipboard \NC \NR 
\NC Control-Delete \NC selection to clipboard                   \NC \NR 
\stoptabulate

When you know the menu items and their locations, you can
use the control key to quicky select the item. 

\starttabulate
\NC Control-<n> \NC menu item <n> \NC \NR 
\stoptabulate

The \type {escape} key is indeed an escape. 

\starttabulate
\NC Escape       \NC goto file menus  \NC \NR 
\NC Shift-Escape \NC goto parent file \NC \NR 
\stoptabulate

The next few bindings are related to \TEX\ specific editing
and running \TEX\ and \METAPOST. The help key pops up the \type
{showtex} window. There you can see what the command asked
for, i.e. where the cursor is located.

In \TEX\ mode, you can run a whole file or juts the
selection. In the latter case, all commands preceding the
start of the text, are included too. In \METAPOST\ mode, the whole
file is processed, but the figure that has the focus is
displayed. 

When you include a file, and this file is present, you can goto that 
file by clicking (in the text) on the name of the file.
Filenames are shown in black and underlined. 

\starttabulate
\NC Control-h \NC provide help         \NC \NR 
\NC Control-r \NC run program          \NC \NR 
\NC Control-o \NC open file            \NC \NR 
\NC Control-f \NC global find (search) \NC \NR 
\stoptabulate

\subject {Menus}


\subject {Installation} 

environment files 

log window

trace mode 

texstep for testing 

\ColofonPage

\stoptext 
